\documentclass[10pt,a4paper]{amsart}
\usepackage{amsmath,amssymb,amsthm,graphicx}
% Retained minimal preamble for the Stacks Project chapter more-morphisms.tex (More on Morphisms).
% Verbatim excerpt of the chapter preamble: Xy-pic setup (native preamble lines 14-18) and the native
% theorem declarations with the mathematical macros (native preamble lines 154-217). Copied mechanically;
% nothing added, removed or redefined. The wrapper class and the amsmath/amssymb/amsthm/graphicx base
% packages are supplied by the assembler, so the class/externaldocument part of the native preamble is omitted.
% For commutative diagrams we use Xy-pic
\usepackage[all]{xy}

% We use 2cell for 2-commutative diagrams.
\xyoption{2cell}
% Theorem environments.
%
\theoremstyle{plain}
\newtheorem{theorem}[subsection]{Theorem}
\newtheorem{proposition}[subsection]{Proposition}
\newtheorem{lemma}[subsection]{Lemma}

\theoremstyle{definition}
\newtheorem{definition}[subsection]{Definition}
\newtheorem{example}[subsection]{Example}
\newtheorem{exercise}[subsection]{Exercise}
\newtheorem{situation}[subsection]{Situation}

\theoremstyle{remark}
\newtheorem{remark}[subsection]{Remark}
\newtheorem{remarks}[subsection]{Remarks}

\numberwithin{equation}{subsection}

% Macros
%
\def\lim{\mathop{\mathrm{lim}}\nolimits}
\def\colim{\mathop{\mathrm{colim}}\nolimits}
\def\Spec{\mathop{\mathrm{Spec}}}
\def\Hom{\mathop{\mathrm{Hom}}\nolimits}
\def\Ext{\mathop{\mathrm{Ext}}\nolimits}
\def\SheafHom{\mathop{\mathcal{H}\!\mathit{om}}\nolimits}
\def\SheafExt{\mathop{\mathcal{E}\!\mathit{xt}}\nolimits}
\def\Sch{\mathit{Sch}}
\def\Mor{\mathop{\mathrm{Mor}}\nolimits}
\def\Ob{\mathop{\mathrm{Ob}}\nolimits}
\def\Sh{\mathop{\mathit{Sh}}\nolimits}
\def\NL{\mathop{N\!L}\nolimits}
\def\CH{\mathop{\mathrm{CH}}\nolimits}
\def\proetale{{pro\text{-}\acute{e}tale}}
\def\etale{{\acute{e}tale}}
\def\QCoh{\mathit{QCoh}}
\def\Ker{\mathop{\mathrm{Ker}}}
\def\Im{\mathop{\mathrm{Im}}}
\def\Coker{\mathop{\mathrm{Coker}}}
\def\Coim{\mathop{\mathrm{Coim}}}

% Boxtimes
%
\DeclareMathSymbol{\boxtimes}{\mathbin}{AMSa}{"02}

%
% Macros for moduli stacks/spaces
%
\def\QCohstack{\mathcal{QC}\!\mathit{oh}}
\def\Cohstack{\mathcal{C}\!\mathit{oh}}
\def\Spacesstack{\mathcal{S}\!\mathit{paces}}
\def\Quotfunctor{\mathrm{Quot}}
\def\Hilbfunctor{\mathrm{Hilb}}
\def\Curvesstack{\mathcal{C}\!\mathit{urves}}
\def\Polarizedstack{\mathcal{P}\!\mathit{olarized}}
\def\Complexesstack{\mathcal{C}\!\mathit{omplexes}}
% \Pic is the operator that assigns to X its picard group, usage \Pic(X)
% \Picardstack_{X/B} denotes the Picard stack of X over B
% \Picardfunctor_{X/B} denotes the Picard functor of X over B
\def\Pic{\mathop{\mathrm{Pic}}\nolimits}
\def\Picardstack{\mathcal{P}\!\mathit{ic}}
\def\Picardfunctor{\mathrm{Pic}}
\def\Deformationcategory{\mathcal{D}\!\mathit{ef}}
\usepackage{cleveref}
\crefname{originallemma}{original lemma}{original lemmas}
\Crefname{originallemma}{Original lemma}{Original lemmas}
\makeatletter
\newlabel{lemma-thicken-property-morphisms}{{09ZV}{source}{}{sourcepointer-1}{}}
\newlabel{lemma-thicken-property-morphisms@cref}{{[originallemma][0][]09ZV}{[1][0][]source}{}{}{}}
\makeatother
\crefname{originallemma}{original lemma}{original lemmas}
\Crefname{originallemma}{Original lemma}{Original lemmas}
\makeatletter
\newlabel{algebra-lemma-NAK}{{00DV}{source}{}{sourcepointer-2}{}}
\newlabel{algebra-lemma-NAK@cref}{{[originallemma][0][]00DV}{[1][0][]source}{}{}{}}
\makeatother
\crefname{originallemma}{original lemma}{original lemmas}
\Crefname{originallemma}{Original lemma}{Original lemmas}
\makeatletter
\newlabel{algebra-lemma-differentials-base-change}{{00RV}{source}{}{sourcepointer-3}{}}
\newlabel{algebra-lemma-differentials-base-change@cref}{{[originallemma][0][]00RV}{[1][0][]source}{}{}{}}
\makeatother
\crefname{originallemma}{original lemma}{original lemmas}
\Crefname{originallemma}{Original lemma}{Original lemmas}
\makeatletter
\newlabel{morphisms-lemma-base-change-finite-type}{{01T4}{source}{}{sourcepointer-4}{}}
\newlabel{morphisms-lemma-base-change-finite-type@cref}{{[originallemma][0][]01T4}{[1][0][]source}{}{}{}}
\makeatother
\crefname{originallemma}{original lemma}{original lemmas}
\Crefname{originallemma}{Original lemma}{Original lemmas}
\makeatletter
\newlabel{morphisms-lemma-quasi-finite-at-point-characterize}{{01TH}{source}{}{sourcepointer-5}{}}
\newlabel{morphisms-lemma-quasi-finite-at-point-characterize@cref}{{[originallemma][0][]01TH}{[1][0][]source}{}{}{}}
\makeatother
\crefname{originallemma}{original lemma}{original lemmas}
\Crefname{originallemma}{Original lemma}{Original lemmas}
\makeatletter
\newlabel{morphisms-lemma-base-change-quasi-finite}{{01TM}{source}{}{sourcepointer-6}{}}
\newlabel{morphisms-lemma-base-change-quasi-finite@cref}{{[originallemma][0][]01TM}{[1][0][]source}{}{}{}}
\makeatother
\crefname{originallemma}{original lemma}{original lemmas}
\Crefname{originallemma}{Original lemma}{Original lemmas}
\makeatletter
\newlabel{morphisms-lemma-base-change-differentials}{{01V0}{source}{}{sourcepointer-7}{}}
\newlabel{morphisms-lemma-base-change-differentials@cref}{{[originallemma][0][]01V0}{[1][0][]source}{}{}{}}
\makeatother
\crefname{originallemma}{original lemma}{original lemmas}
\Crefname{originallemma}{Original lemma}{Original lemmas}
\makeatletter
\newlabel{morphisms-lemma-base-change-proper}{{01W4}{source}{}{sourcepointer-8}{}}
\newlabel{morphisms-lemma-base-change-proper@cref}{{[originallemma][0][]01W4}{[1][0][]source}{}{}{}}
\makeatother
\crefname{originallemma}{original lemma}{original lemmas}
\Crefname{originallemma}{Original lemma}{Original lemmas}
\makeatletter
\newlabel{morphisms-lemma-finite-integral}{{01WJ}{source}{}{sourcepointer-9}{}}
\newlabel{morphisms-lemma-finite-integral@cref}{{[originallemma][0][]01WJ}{[1][0][]source}{}{}{}}
\makeatother
\crefname{originallemma}{original lemma}{original lemmas}
\Crefname{originallemma}{Original lemma}{Original lemmas}
\makeatletter
\newlabel{morphisms-lemma-base-change-finite}{{01WL}{source}{}{sourcepointer-10}{}}
\newlabel{morphisms-lemma-base-change-finite@cref}{{[originallemma][0][]01WL}{[1][0][]source}{}{}{}}
\makeatother
\crefname{originallemma}{original lemma}{original lemmas}
\Crefname{originallemma}{Original lemma}{Original lemmas}
\makeatletter
\newlabel{schemes-lemma-base-change-immersion}{{01JY}{source}{}{sourcepointer-11}{}}
\newlabel{schemes-lemma-base-change-immersion@cref}{{[originallemma][0][]01JY}{[1][0][]source}{}{}{}}
\makeatother
\crefname{originallemma}{original lemma}{original lemmas}
\Crefname{originallemma}{Original lemma}{Original lemmas}
\makeatletter
\newlabel{morphisms-lemma-base-change-unramified}{{02GA}{source}{}{sourcepointer-12}{}}
\newlabel{morphisms-lemma-base-change-unramified@cref}{{[originallemma][0][]02GA}{[1][0][]source}{}{}{}}
\makeatother
\crefname{originallemma}{original lemma}{original lemmas}
\Crefname{originallemma}{Original lemma}{Original lemmas}
\makeatletter
\newlabel{morphisms-lemma-base-change-relative-dimension-d}{{02NK}{source}{}{sourcepointer-13}{}}
\newlabel{morphisms-lemma-base-change-relative-dimension-d@cref}{{[originallemma][0][]02NK}{[1][0][]source}{}{}{}}
\makeatother
\crefname{originaldefinition}{original definition}{original definitions}
\Crefname{originaldefinition}{Original definition}{Original definitions}
\makeatletter
\newlabel{morphisms-definition-relative-dimension-d}{{02NJ}{source}{}{sourcepointer-14}{}}
\newlabel{morphisms-definition-relative-dimension-d@cref}{{[originaldefinition][0][]02NJ}{[1][0][]source}{}{}{}}
\makeatother
\crefname{originallemma}{original lemma}{original lemmas}
\Crefname{originallemma}{Original lemma}{Original lemmas}
\makeatletter
\newlabel{schemes-lemma-base-change-monomorphism}{{02YC}{source}{}{sourcepointer-15}{}}
\newlabel{schemes-lemma-base-change-monomorphism@cref}{{[originallemma][0][]02YC}{[1][0][]source}{}{}{}}
\makeatother
\crefname{originallemma}{original lemma}{original lemmas}
\Crefname{originallemma}{Original lemma}{Original lemmas}
\makeatletter
\newlabel{morphisms-lemma-differentials-diagonal}{{08S2}{source}{}{sourcepointer-16}{}}
\newlabel{morphisms-lemma-differentials-diagonal@cref}{{[originallemma][0][]08S2}{[1][0][]source}{}{}{}}
\makeatother
\title{Properties of morphisms preserved by a cartesian finite order thickening of the base}
\author{The Stacks Project Authors}
\date{stacks-project:more-morphisms.tex:local-snapshot-20261011; GNU FDL 1.2 or later}
\begin{document}
\maketitle
\noindent\textit{Editorial note.} All delivered mathematical text is the author's own, copied from the pinned Stacks Project chapter file more-morphisms.tex; no statement or proof step was rewritten, reordered or removed, including the author's own eleventh item 'add more here'.\par\medskip
\noindent\textit{Editorial note.} The two leading context spans are the chapter's own definition of thickenings, first order thickenings, finite order thickenings and morphisms of thickenings (label definition-thickening), and the author's discussion of the filtration of an nth order thickening by first order thickenings that the proof cites at line 379; both define or explain notions appearing in the counted statement, and definition-thickening is therefore not listed among the external references.\par\medskip
\noindent\textit{Editorial note.} The preamble is the chapter's own retained minimal preamble (sources/stacks-more-morphisms/retained-preamble.tex): a verbatim excerpt of the native Xy-pic setup and of the native theorem declarations and mathematical macros. The native class line, the xr-hyper externaldocument block and the reference/slogan/history display environments are omitted because the assembler supplies a plain amsart wrapper and the excerpt uses none of those environments; no macro was invented or redefined.\par\medskip
\noindent\textit{Editorial note.} Locators of the cross-chapter pointers refer to a verbatim excerpt of the project's own tag list, sources/stacks-more-morphisms/referenced-tags.txt, each line carrying its stable Stacks tag; the same-chapter pointer to lemma-thicken-property-morphisms refers to its own label line inside the pinned chapter file.\par\medskip
\noindent\textit{Allowed standard prerequisites (editorial).}\begin{itemize}
\item Stability of the listed morphism properties under base change. Schemes, Stacks tags 01JY and 02YC, and Morphisms, Stacks tags 01T4, 01TM, 02NK, 01V0, 02GA, 01W4 and 01WL, all displayed at proof lines 358-371. Gives the easy direction of every clause: if f' has the property then so does its base change f along X = S times\_\{S'\} X'.
\item Nakayama's lemma. Algebra, Stacks tag 00DV (algebra-lemma-NAK), invoked at proof lines 405, 410 and 430. Turns surjectivity of A to B into surjectivity of A' to B', lifts finite typeness along A'[x\_1,...,x\_n] to B', and upgrades Omega\_\{B'/A'\} = I Omega\_\{B'/A'\} to Omega\_\{B'/A'\} = 0.
\item Reduction of a finite order thickening to first order thickenings. the chapter's own discussion printed as the second context span (source lines 66-85), cited by the proof at line 379. Justifies the induction on the order of the thickening with which the proof may assume that S in S' is a first order thickening.
\item Permanence of quasi-compactness, separatedness and universal closedness under thickenings. More on Morphisms, Stacks tag 09ZV (more-morphisms-lemma-thicken-property-morphisms), invoked at proof lines 441 and 447. Supplies the remaining ingredients for clauses (7), proper, and (8), finite, once finite typeness and the integral or unramified part is available.
\item Fibrewise criteria for quasi-finiteness and for relative dimension d, and the diagonal criterion for differentials. Morphisms, Stacks tags 01TH, 02NJ (morphisms-definition-relative-dimension-d) and 08S2, invoked at proof lines 415-422 and 462-464. Lets clauses (3) and (4) be checked on fibres, and lets clause (9) reduce the monomorphism statement to the vanishing of Omega\_\{X'/S'\}.
\end{itemize}\medskip
\ifdefined\hypertarget\hypertarget{sourcepointer-1}{}\fi
\noindent\textit{Source reference.} Original lemma 09ZV: sources/stacks-more-morphisms/pinned-original-source.tex, source lines 270--270. Same-chapter input invoked at proof lines 441 and 447 for clauses (7) and (8); a permanence statement about quasi-compactness, separatedness and universal closedness that the counted proof combines with the finite type clauses it proves itself.\par\medskip
\ifdefined\hypertarget\hypertarget{sourcepointer-2}{}\fi
\noindent\textit{Source reference.} Original lemma 00DV: sources/stacks-more-morphisms/referenced-tags.txt, source lines 4--4. Nakayama's lemma invoked three times, at proof lines 405, 410 and 430; a standard algebraic input, while the affine reduction that makes it applicable is part of the counted proof.\par\medskip
\ifdefined\hypertarget\hypertarget{sourcepointer-3}{}\fi
\noindent\textit{Source reference.} Original lemma 00RV: sources/stacks-more-morphisms/referenced-tags.txt, source lines 5--5. Identification of Omega\_\{B/A\} with Omega\_\{B'/A'\} modulo I Omega\_\{B'/A'\} at proof line 427; a navigation pointer, the Nakayama step that consumes it being in the counted proof.\par\medskip
\ifdefined\hypertarget\hypertarget{sourcepointer-4}{}\fi
\noindent\textit{Source reference.} Original lemma 01T4: sources/stacks-more-morphisms/referenced-tags.txt, source lines 6--6. Base change stability of local finite typeness, displayed at proof lines 364-365 as part of the uniform easy direction.\par\medskip
\ifdefined\hypertarget\hypertarget{sourcepointer-5}{}\fi
\noindent\textit{Source reference.} Original lemma 01TH: sources/stacks-more-morphisms/referenced-tags.txt, source lines 7--7. Fibrewise criterion for quasi-finiteness used at proof line 417 to deduce clause (3) from clause (2).\par\medskip
\ifdefined\hypertarget\hypertarget{sourcepointer-6}{}\fi
\noindent\textit{Source reference.} Original lemma 01TM: sources/stacks-more-morphisms/referenced-tags.txt, source lines 8--8. Base change stability of local quasi-finiteness, displayed at proof line 366 as part of the uniform easy direction.\par\medskip
\ifdefined\hypertarget\hypertarget{sourcepointer-7}{}\fi
\noindent\textit{Source reference.} Original lemma 01V0: sources/stacks-more-morphisms/referenced-tags.txt, source lines 9--9. Base change stability of the vanishing of differentials, displayed at proof line 368; the algebraic translation of clause (5) is in the counted proof.\par\medskip
\ifdefined\hypertarget\hypertarget{sourcepointer-8}{}\fi
\noindent\textit{Source reference.} Original lemma 01W4: sources/stacks-more-morphisms/referenced-tags.txt, source lines 10--10. Base change stability of properness, displayed at proof line 370 as part of the uniform easy direction.\par\medskip
\ifdefined\hypertarget\hypertarget{sourcepointer-9}{}\fi
\noindent\textit{Source reference.} Original lemma 01WJ: sources/stacks-more-morphisms/referenced-tags.txt, source lines 11--11. Description of finite morphisms as integral and locally of finite type, used at proof line 449 for clause (8).\par\medskip
\ifdefined\hypertarget\hypertarget{sourcepointer-10}{}\fi
\noindent\textit{Source reference.} Original lemma 01WL: sources/stacks-more-morphisms/referenced-tags.txt, source lines 12--12. Base change stability of finiteness, displayed at proof line 371 as part of the uniform easy direction.\par\medskip
\ifdefined\hypertarget\hypertarget{sourcepointer-11}{}\fi
\noindent\textit{Source reference.} Original lemma 01JY: sources/stacks-more-morphisms/referenced-tags.txt, source lines 14--14. Base change stability of immersions, displayed at proof line 361 as part of the uniform easy direction for clause (10).\par\medskip
\ifdefined\hypertarget\hypertarget{sourcepointer-12}{}\fi
\noindent\textit{Source reference.} Original lemma 02GA: sources/stacks-more-morphisms/referenced-tags.txt, source lines 15--15. Base change stability of unramifiedness, displayed at proof line 369 as part of the uniform easy direction.\par\medskip
\ifdefined\hypertarget\hypertarget{sourcepointer-13}{}\fi
\noindent\textit{Source reference.} Original lemma 02NK: sources/stacks-more-morphisms/referenced-tags.txt, source lines 16--16. Base change stability of the relative dimension, displayed at proof line 367; clause (4) is then checked on fibres inside the counted proof.\par\medskip
\ifdefined\hypertarget\hypertarget{sourcepointer-14}{}\fi
\noindent\textit{Source reference.} Original definition 02NJ: sources/stacks-more-morphisms/referenced-tags.txt, source lines 17--17. Definition of locally of finite type of relative dimension d, cited at proof line 422 for the fibrewise reading of clause (4).\par\medskip
\ifdefined\hypertarget\hypertarget{sourcepointer-15}{}\fi
\noindent\textit{Source reference.} Original lemma 02YC: sources/stacks-more-morphisms/referenced-tags.txt, source lines 18--18. Base change stability of monomorphisms, displayed at proof line 362 as part of the uniform easy direction for clause (9).\par\medskip
\ifdefined\hypertarget\hypertarget{sourcepointer-16}{}\fi
\noindent\textit{Source reference.} Original lemma 08S2: sources/stacks-more-morphisms/referenced-tags.txt, source lines 20--20. Diagonal criterion for differentials invoked at proof line 462 to finish clause (9) from the vanishing of Omega\_\{X'/S'\}; the surrounding computation of the ideal of the diagonal is in the counted proof.\par\medskip
\begin{definition}
\label{definition-thickening}
Thickenings.
\begin{enumerate}
\item We say a scheme $X'$ is a {\it thickening} of a scheme $X$ if
$X$ is a closed subscheme of $X'$ and the underlying topological spaces
are equal.
\item We say a scheme $X'$ is a {\it first order thickening} of a scheme $X$ if
$X$ is a closed subscheme of $X'$ and the quasi-coherent sheaf of ideals
$\mathcal{I} \subset \mathcal{O}_{X'}$ defining $X$ has square zero.
\item We say a scheme $X'$ is a {\it finite order thickening} of a scheme $X$
if $X$ is a closed subscheme of $X'$ and the quasi-coherent sheaf of ideals
$\mathcal{I} \subset \mathcal{O}_{X'}$ defining $X$ is nilpotent, i.e.,
there exists an integer $n \geq 0$ such that $\mathcal{I}^{n + 1} = 0$.
\item We say a scheme $X'$ is an {\it $n$th order thickening} of a scheme $X$
if $X$ is a closed subscheme of $X$ and $\mathcal{I}^{n + 1} = 0$
where $\mathcal{I} \subset \mathcal{O}_{X'}$ is the
quasi-coherent sheaf of ideals defining $X$.
\item Given two thickenings $X \subset X'$ and $Y \subset Y'$ a
{\it morphism of thickenings} is a morphism $f' : X' \to Y'$ such that
$f'(X) \subset Y$, i.e., such that $f'|_X$ factors through the closed
subscheme $Y$. In this situation we set $f = f'|_X : X \to Y$ and we say
that $(f, f') : (X \subset X') \to (Y \subset Y')$ is a morphism of
thickenings.
\item Let $S$ be a scheme. We similarly define {\it thickenings over $S$}, and
{\it morphisms of thickenings over $S$}. This means that the schemes
$X, X', Y, Y'$ above are schemes over $S$, and that the morphisms
$X \to X'$, $Y \to Y'$ and $f' : X' \to Y'$ are morphisms over $S$.
\end{enumerate}
\end{definition}

\noindent
Let $i_X : X \to X'$ be a thickening. Any local section of the kernel
$\mathcal{I} = \Ker(i_X^\sharp)$ is locally nilpotent. This gives
some control over the structure sheaf of $\mathcal{O}_{X'}$, but
technically the class of finite order thickenings $X \subset X'$
is much easier to handle. Namely, if $X'$ is an $n$th order thinking
then we have a filtration
$$
0 = \mathcal{I}^{n + 1} \subset
\mathcal{I}^n \subset
\mathcal{I}^{n - 1} \subset \ldots \subset
\mathcal{I} \subset \mathcal{O}_{X'}
$$
and we see that $X'$ is filtered by closed subspaces
$$
X = X_1 \subset X_2 \subset \ldots \subset X_n \subset X_{n + 1} = X'
$$
such that each pair $X_i \subset X_{i + 1}$ is a first order thickening
over $S$. Using simple induction arguments many results proved for first order
thickenings can be rephrased as results on finite order thickenings.

\begin{lemma}
\label{lemma-thicken-property-morphisms-cartesian}
Let $(f, f') : (X \subset X') \to (S \subset S')$ be a morphism
of thickenings such that $X = S \times_{S'} X'$. If $S \subset S'$
is a finite order thickening, then
\begin{enumerate}
\item $f$ is a closed immersion if and only if $f'$ is a closed immersion,
\item $f$ is locally of finite type if and only if $f'$ is
locally of finite type,
\item $f$ is locally quasi-finite if and only if $f'$ is locally
quasi-finite,
\item $f$ is locally of finite type of relative dimension $d$ if and
only if $f'$ is locally of finite type of relative dimension $d$,
\item $\Omega_{X/S} = 0$ if and only if $\Omega_{X'/S'} = 0$,
\item $f$ is unramified if and only if $f'$ is unramified,
\item $f$ is proper if and only if $f'$ is proper,
\item $f$ is finite if and only if $f'$ is finite,
\item $f$ is a monomorphism if and only if $f'$ is a monomorphism,
\item $f$ is an immersion if and only if $f'$ is an immersion, and
\item add more here.
\end{enumerate}
\end{lemma}

\begin{proof}
The properties $\mathcal{P}$ listed in the lemma are all stable
under base change, hence if $f'$ has property $\mathcal{P}$, then so
does $f$. See
Schemes, Lemmas \ref{schemes-lemma-base-change-immersion} and
\ref{schemes-lemma-base-change-monomorphism}
and
Morphisms, Lemmas
\ref{morphisms-lemma-base-change-finite-type},
\ref{morphisms-lemma-base-change-quasi-finite},
\ref{morphisms-lemma-base-change-relative-dimension-d},
\ref{morphisms-lemma-base-change-differentials},
\ref{morphisms-lemma-base-change-unramified},
\ref{morphisms-lemma-base-change-proper}, and
\ref{morphisms-lemma-base-change-finite}.

\medskip\noindent
The interesting direction in each case is therefore to assume
that $f$ has the property and deduce that $f'$ has it too.
By induction on the order of the thickening we may
assume that $S \subset S'$ is a first order thickening, see
discussion immediately following
Definition \ref{definition-thickening}.

\medskip\noindent
Most of the proofs will use a reduction to the affine case. Let
$U' \subset S'$ be an affine open and let $V' \subset X'$ be an affine open
lying over $U'$. Let $U' = \Spec(A')$ and denote $I \subset A'$ be the ideal
defining the closed subscheme $U' \cap S$. Say $V' = \Spec(B')$.
Then $V' \cap X = \Spec(B'/IB')$. Setting $A = A'/I$ and
$B = B'/IB'$ we get a commutative diagram
$$
\xymatrix{
0 \ar[r] &
IB' \ar[r] &
B' \ar[r] &
B \ar[r] & 0 \\
0 \ar[r] &
IA' \ar[r] \ar[u] &
A' \ar[r] \ar[u] &
A \ar[r] \ar[u] & 0
}
$$
with exact rows and $I^2 = 0$.

\medskip\noindent
The translation of (1) into algebra: If $A \to B$ is surjective,
then $A' \to B'$ is surjective. This follows from
Nakayama's lemma (Algebra, Lemma \ref{algebra-lemma-NAK}).

\medskip\noindent
The translation of (2) into algebra: If $A \to B$ is a finite type ring
map, then $A' \to B'$ is a finite type ring map. This follows from
Nakayama's lemma (Algebra, Lemma \ref{algebra-lemma-NAK})
applied to a map $A'[x_1, \ldots, x_n] \to B'$ such that
$A[x_1, \ldots, x_n] \to B$ is surjective.

\medskip\noindent
Proof of (3).  Follows from (2) and that quasi-finiteness of a morphism
which is locally of finite type can be checked on fibres, see
Morphisms, Lemma \ref{morphisms-lemma-quasi-finite-at-point-characterize}.

\medskip\noindent
Proof of (4). Follows from (2) and that the additional property of ``being of
relative dimension $d$'' can be checked on fibres (by definition, see
Morphisms, Definition \ref{morphisms-definition-relative-dimension-d}.

\medskip\noindent
The translation of (5) into algebra: If $\Omega_{B/A} = 0$, then
$\Omega_{B'/A'} = 0$. By
Algebra, Lemma \ref{algebra-lemma-differentials-base-change}
we have $0 = \Omega_{B/A} = \Omega_{B'/A'}/I\Omega_{B'/A'}$.
Hence $\Omega_{B'/A'} = 0$ by
Nakayama's lemma (Algebra, Lemma \ref{algebra-lemma-NAK}).

\medskip\noindent
The translation of (6) into algebra: If $A \to B$ is unramified
map, then $A' \to B'$ is unramified. Since $A \to B$ is of finite
type we see that $A' \to B'$ is of finite type by (2) above.
Since $A \to B$ is unramified we have $\Omega_{B/A} = 0$. By
part (5) we have $\Omega_{B'/A'} = 0$. Thus $A' \to B'$ is unramified.

\medskip\noindent
Proof of (7). Follows by combining (2) with
results of Lemma \ref{lemma-thicken-property-morphisms}
and the fact that proper equals quasi-compact $+$
separated $+$ locally of finite type $+$ universally closed.

\medskip\noindent
Proof of (8). Follows by combining (2) with
results of Lemma \ref{lemma-thicken-property-morphisms}
and using the fact that finite equals integral $+$ locally
of finite type (Morphisms, Lemma \ref{morphisms-lemma-finite-integral}).

\medskip\noindent
Proof of (9). As $f$ is a monomorphism we have $X = X \times_S X$.
We may apply the results proved so far to the morphism of thickenings
$(X \subset X') \to (X \times_S X \subset X' \times_{S'} X')$.
We conclude $X' \to X' \times_{S'} X'$ is a closed immersion by (1).
In fact, it is a first order thickening as the ideal defining the
closed immersion
$X' \to X' \times_{S'} X'$ is contained in the pullback of the ideal
$\mathcal{I} \subset \mathcal{O}_{S'}$ cutting out $S$ in $S'$.
Indeed, $X = X \times_S X = (X' \times_{S'} X') \times_{S'} S$
is contained in $X'$. Hence by
Morphisms, Lemma \ref{morphisms-lemma-differentials-diagonal}
it suffices to show that
$\Omega_{X'/S'} = 0$ which follows from (5)
and the corresponding statement for $X/S$.

\medskip\noindent
Proof of (10). If $f : X \to S$ is an immersion, then it factors as
$X \to U \to S$ where $U \to S$ is an open immersion and $X \to U$ is a
closed immersion. Let $U' \subset S'$ be the open subscheme whose
underlying topological space is the same as $U$. Then $X' \to S'$
factors through $U'$ and we conclude that $X' \to U'$ is a closed
immersion by part (1). This finishes the proof.
\end{proof}

\end{document}
